% GENERATED FILE. Edit canonical lessons, then run npm run generate:books.
% canonical-source-hash: fnv1a64:1a7181b15ae3546e
% canonical-lessons: HI-C193-adhyapak, HI-C193-vakil, HI-C193-sipahi, HI-C193-darzi, HI-C193-nai

\chapter{A teacher, A lawyer, A soldier, A tailor, A barber}
\label{ch:hi-a-teacher-a-lawyer-a-soldier-a-tailor-a-barber}
\hlchaptermodality{193}

\begin{chapteropening}
\textbf{By the end of this chapter:} \emph{I can say a teacher, a lawyer, a soldier, a tailor and a barber.}

Complete the last lesson of chapter 193: I can say a teacher, a lawyer, a soldier, a tailor and a barber.
\end{chapteropening}

\section[\texorpdfstring{\hi{अध्यापक} (adhyāpak)}{अध्यापक (adhyāpak)}]{\hi{अध्यापक} (adhyāpak) — a teacher}

\label{lesson:HI-C193-adhyapak}



\begin{quote}
Before the new one: say the Hindi for often, then the Hindi for never.
\end{quote}

\subsection*{You'll want to know: \hi{अध्यापक}}
\textbf{\hi{अध्यापक}} — \emph{adhyāpak} — \textquotedblleft{}a teacher\textquotedblright{}.

\begin{cousinweb}
\textbf{\hi{अध्यापक}} (\emph{adhyāpak}) is Sanskrit \emph{adhyāpaka}, \textquotedblleft{}teacher\textquotedblright{}, taken over whole as a learned loan, a tatsama (CDIAL 278).
\end{cousinweb}

\begin{grammarlens}[title={What you've built}]
One more word for this chapter.
\end{grammarlens}

\subsection*{Your turn}
\begin{itemize}
\raggedright
  \item \emph{Say it:} \emph{adhyāpak}
  \item \emph{Say it:} \emph{adhyāpak}, once more
  \item \emph{Say it:} say \emph{adhyāpak}
  \item \emph{Recall:} say the Hindi for often, then the Hindi for never, then say \emph{adhyāpak} again
\end{itemize}

\begin{tcolorbox}[breakable,colback=teal!4,colframe=teal!35!black,title={Before you move on}]
What does \hi{अध्यापक} mean? (\textquotedblleft{}A teacher\textquotedblright{}.) Say it once more.
\end{tcolorbox}
\section[\texorpdfstring{\hi{वकील} (vakīl)}{वकील (vakīl)}]{\hi{वकील} (vakīl) — a lawyer}

\label{lesson:HI-C193-vakil}



\begin{quote}
Before the new one: say the Hindi for never, then the Hindi for a teacher.
\end{quote}

\subsection*{You'll want to know: \hi{वकील}}
\textbf{\hi{वकील}} — \emph{vakīl} — \textquotedblleft{}a lawyer\textquotedblright{}.

\begin{cousinweb}
\textbf{\hi{वकील}} (\emph{vakīl}) is a loan from Arabic \emph{wakīl}, \textquotedblleft{}an agent, a representative\textquotedblright{}, which reached Hindi through Persian. Its Arabic root is \emph{w-k-l} (Lane).
\end{cousinweb}

\begin{grammarlens}[title={What you've built}]
One more word for this chapter.
\end{grammarlens}

\subsection*{Your turn}
\begin{itemize}
\raggedright
  \item \emph{Say it:} \emph{vakīl}
  \item \emph{Say it:} \emph{vakīl}, once more
  \item \emph{Say it:} say \emph{vakīl}
  \item \emph{Recall:} say the Hindi for never, then the Hindi for a teacher, then say \emph{vakīl} again
\end{itemize}

\begin{tcolorbox}[breakable,colback=teal!4,colframe=teal!35!black,title={Before you move on}]
What does \hi{वकील} mean? (\textquotedblleft{}A lawyer\textquotedblright{}.) Say it once more.
\end{tcolorbox}
\section[\texorpdfstring{\hi{सिपाही} (sipāhī)}{सिपाही (sipāhī)}]{\hi{सिपाही} (sipāhī) — a soldier, a constable}

\label{lesson:HI-C193-sipahi}



\begin{quote}
Before the new one: say the Hindi for a teacher, then the Hindi for a lawyer.
\end{quote}

\subsection*{You'll want to know: \hi{सिपाही}}
\textbf{\hi{सिपाही}} — \emph{sipāhī} — \textquotedblleft{}a soldier, a constable\textquotedblright{}.

\begin{cousinweb}
\textbf{\hi{सिपाही}} (\emph{sipāhī}) is a loan, usually traced to Persian \emph{sipāhī}, \textquotedblleft{}soldier\textquotedblright{}. It is built on \emph{sipāh}, \textquotedblleft{}army\textquotedblright{}. English took the word as \emph{sepoy}.
\end{cousinweb}

\begin{grammarlens}[title={What you've built}]
One more word for this chapter.
\end{grammarlens}

\subsection*{Your turn}
\begin{itemize}
\raggedright
  \item \emph{Say it:} \emph{sipāhī}
  \item \emph{Say it:} \emph{sipāhī}, once more
  \item \emph{Say it:} say \emph{sipāhī}
  \item \emph{Recall:} say the Hindi for a teacher, then the Hindi for a lawyer, then say \emph{sipāhī} again
\end{itemize}

\begin{tcolorbox}[breakable,colback=teal!4,colframe=teal!35!black,title={Before you move on}]
What does \hi{सिपाही} mean? (\textquotedblleft{}A soldier, a constable\textquotedblright{}.) Say it once more.
\end{tcolorbox}
\section[\texorpdfstring{\hi{दर्ज़ी} (darzī)}{दर्ज़ी (darzī)}]{\hi{दर्ज़ी} (darzī) — a tailor}

\label{lesson:HI-C193-darzi}



\begin{quote}
Before the new one: say the Hindi for a lawyer, then the Hindi for a soldier.
\end{quote}

\subsection*{You'll want to know: \hi{दर्ज़ी}}
\textbf{\hi{दर्ज़ी}} — \emph{darzī} — \textquotedblleft{}a tailor\textquotedblright{}.

\begin{cousinweb}
\textbf{\hi{दर्ज़ी}} (\emph{darzī}) is a loan, usually traced to Persian \emph{darzī}, \textquotedblleft{}tailor\textquotedblright{}. It is built on \emph{darz}, \textquotedblleft{}seam\textquotedblright{}.
\end{cousinweb}

\begin{grammarlens}[title={What you've built}]
One more word for this chapter.
\end{grammarlens}

\subsection*{Your turn}
\begin{itemize}
\raggedright
  \item \emph{Say it:} \emph{darzī}
  \item \emph{Say it:} \emph{darzī}, once more
  \item \emph{Say it:} say \emph{darzī}
  \item \emph{Recall:} say the Hindi for a lawyer, then the Hindi for a soldier, then say \emph{darzī} again
\end{itemize}

\begin{tcolorbox}[breakable,colback=teal!4,colframe=teal!35!black,title={Before you move on}]
What does \hi{दर्ज़ी} mean? (\textquotedblleft{}A tailor\textquotedblright{}.) Say it once more.
\end{tcolorbox}
\section[\texorpdfstring{\hi{नाई} (nāī)}{नाई (nāī)}]{\hi{नाई} (nāī) — a barber}

\label{lesson:HI-C193-nai}



\begin{quote}
Before the new one: say the Hindi for a soldier, then the Hindi for a tailor.
\end{quote}

\subsection*{You'll want to know: \hi{नाई}}
\textbf{\hi{नाई}} — \emph{nāī} — \textquotedblleft{}a barber\textquotedblright{}.

\begin{cousinweb}
\textbf{\hi{नाई}} (\emph{nāī}) is a tadbhava, inherited and reshaped by sound change. CDIAL lists it, as a dialect form, under Sanskrit \emph{nāpita}, \textquotedblleft{}barber\textquotedblright{}. The by-form \emph{nāū}, \textquotedblleft{}barber\textquotedblright{}, sits in the same entry under a second word, \emph{snāpaka}, \textquotedblleft{}bath attendant\textquotedblright{}, built on \emph{snā-}, \textquotedblleft{}to bathe\textquotedblright{} (CDIAL 13792).
\end{cousinweb}

\begin{grammarlens}[title={What you've built}]
One more word for this chapter.
\end{grammarlens}

\subsection*{Your turn}
\begin{itemize}
\raggedright
  \item \emph{Say it:} \emph{nāī}
  \item \emph{Say it:} \emph{nāī}, once more
  \item \emph{Say it:} say \emph{nāī}
  \item \emph{Recall:} say the Hindi for a soldier, then the Hindi for a tailor, then say \emph{nāī} again
\end{itemize}

\begin{tcolorbox}[breakable,colback=teal!4,colframe=teal!35!black,title={Before you move on}]
What does \hi{नाई} mean? (\textquotedblleft{}A barber\textquotedblright{}.) Say it once more.
\end{tcolorbox}
